\chapter{The Mechanical Derivation of Optical Relativity Limits}
\label{app:optical_relativity}
\addtocontents{toc}{\protect\setcounter{tocdepth}{0}}

This appendix provides the explicit weak-field mathematical evaluations for the macroscopic integrations of the strain-induced vacuum refractive index derived in \hyperref[sec:refractive_index]{Section~\ref*{sec:refractive_index}}. In the gravitational optical limit, the active-vacuum refractive index is
\begin{equation}
    n_\tau(r) = \frac{1}{(1-\alpha_{\mathrm{G}})^2}.
\end{equation}

For weak gravitational strain, where $\alpha_{\mathrm{G}} \ll 1$, this reduces to
\begin{equation}
    n_\tau(r) \approx 1+2\alpha_{\mathrm{G}} = 1+\frac{2GM}{r c_{\mathrm{E}}^2}.
\end{equation}

The following integrations demonstrate how this strain-induced refractive index reproduces the standard weak-field optical limits associated with gravitational light deflection and Shapiro time delay.

\section{Evaluation of the Gravitational Lensing Deflection Integral}
As established in Tension Field Theory, the total deflection angle $\theta$ of a photon wave packet is the accumulated transverse refractive-gradient steering along the unperturbed line-of-sight path, taken here as the $z$-axis.

Given the impact parameter $b$, the radial distance from the mass is
\begin{equation}
    r = \sqrt{b^2 + z^2}.
\end{equation}

The deflection magnitude is governed by
\begin{equation}
    \theta = \left| \int_{-\infty}^{+\infty} \nabla_{\perp} n_\tau \,\mathrm{d}z \right| = \left| \int_{-\infty}^{+\infty} \frac{\partial}{\partial b} \left( 1+ \frac{2GM}{c_{\mathrm{E}}^2\sqrt{b^2+z^2}} \right) \mathrm{d}z \right|.
\end{equation}

\subsection{Step 1: Compute the Partial Derivative}
Differentiating the refractive index with respect to the impact parameter $b$ yields the transverse gradient:
\begin{equation}
    \frac{\partial n_\tau}{\partial b} = -\frac{2GMb}{c_{\mathrm{E}}^2(b^2+z^2)^{3/2}}.
\end{equation}
The negative sign indicates inward steering toward the mass. Taking the magnitude gives the total deflection angle:
\begin{equation}
    \theta = \frac{2GMb}{c_{\mathrm{E}}^2} \int_{-\infty}^{+\infty} \frac{\mathrm{d}z}{(b^2+z^2)^{3/2}}.
\end{equation}

\subsection{Step 2: Trigonometric Substitution}
To evaluate the integral over the full path, we apply the substitution
\begin{equation}
    z = b\tan\phi,
\end{equation}
which implies
\begin{equation}
    \mathrm{d}z = b\sec^2\phi\,\mathrm{d}\phi.
\end{equation}
The integration limits change from $z \in (-\infty, +\infty)$ to $\phi \in (-\pi/2, +\pi/2)$.

The denominator becomes
\begin{equation}
    (b^2+b^2\tan^2\phi)^{3/2} = (b^2\sec^2\phi)^{3/2} = b^3\sec^3\phi.
\end{equation}

Substituting these into the integral gives
\begin{equation}
    \int_{-\pi/2}^{+\pi/2} \frac{b\sec^2\phi\,\mathrm{d}\phi}{b^3\sec^3\phi} = \frac{1}{b^2} \int_{-\pi/2}^{+\pi/2} \cos\phi\,\mathrm{d}\phi.
\end{equation}

\subsection{Step 3: Final Integration}
Evaluating the cosine integral:
\begin{equation}
    \frac{1}{b^2} \left[ \sin\phi \right]_{-\pi/2}^{+\pi/2} = \frac{1}{b^2} \left( 1-(-1) \right) = \frac{2}{b^2}.
\end{equation}

Substituting this result back into the deflection equation gives
\begin{equation}
    \theta_{\mathrm{total}} = \frac{2GMb}{c_{\mathrm{E}}^2} \left( \frac{2}{b^2} \right) = \frac{4GM}{bc_{\mathrm{E}}^2}.
\end{equation}

Because the local lattice strain at the impact parameter is defined as
\begin{equation}
    \alpha_{\mathrm{G}}(b) = \frac{GM}{bc_{\mathrm{E}}^2},
\end{equation}
the total macroscopic deflection is
\begin{equation}
    \theta_{\mathrm{total}} = 4\alpha_{\mathrm{G}}(b).
\end{equation}

Thus, the standard weak-field Einstein deflection angle is recovered as the accumulated transverse steering of a photon wave packet through the strain-index gradient of the active vacuum.

\section{Evaluation of the Shapiro Time Delay Integral}
The local propagation speed of a signal is inversely proportional to the strain-induced refractive index:
\begin{equation}
    c_{\mathrm{P}}(r) = \frac{c_{\mathrm{E}}}{n_\tau(r)}.
\end{equation}

In the weak-field limit,
\begin{equation}
    n_\tau(r) \approx 1+2\alpha_{\mathrm{G}},
\end{equation}
so the local propagation speed becomes
\begin{equation}
    c_{\mathrm{P}}(r) \approx c_{\mathrm{E}} \left( 1-2\alpha_{\mathrm{G}} \right).
\end{equation}

The total flight time $\Delta t$ for a signal traveling from an emitter at $-z_1$ to a receiver at $+z_2$ is the integral of the inverse velocity along the geometric path:
\begin{equation}
    \Delta t = \int_{-z_1}^{+z_2} \frac{1}{c_{\mathrm{P}}} \,\mathrm{d}z \approx \frac{1}{c_{\mathrm{E}}} \int_{-z_1}^{+z_2} \left( 1+ \frac{2GM}{c_{\mathrm{E}}^2\sqrt{b^2+z^2}} \right) \mathrm{d}z.
\end{equation}

\subsection{Step 1: Isolate the Delay Component}
The integral separates into the unstrained baseline travel time and the optical delay induced by the refractive vacuum:
\begin{equation}
    \Delta t = \frac{z_1+z_2}{c_{\mathrm{E}}} + \Delta t_{\mathrm{delay}}.
\end{equation}

The delay component is therefore
\begin{equation}
    \Delta t_{\mathrm{delay}} = \frac{2GM}{c_{\mathrm{E}}^3} \int_{-z_1}^{+z_2} \frac{\mathrm{d}z}{\sqrt{b^2+z^2}}.
\end{equation}

\subsection{Step 2: Integration of the Delay Function}
The standard integral of $(b^2+z^2)^{-1/2}$ gives a natural logarithm. Evaluating from the emitter to the receiver yields
\begin{equation}
    \Delta t_{\mathrm{delay}} = \frac{2GM}{c_{\mathrm{E}}^3} \left[ \ln \left( z+\sqrt{b^2+z^2} \right) \right]_{-z_1}^{+z_2}.
\end{equation}
Therefore,
\begin{equation}
    \Delta t_{\mathrm{delay}} = \frac{2GM}{c_{\mathrm{E}}^3} \ln \left( \frac{ z_2+\sqrt{b^2+z_2^2} }{ -z_1+\sqrt{b^2+z_1^2} } \right).
\end{equation}

\subsection{Step 3: Macroscopic Approximation}
For planetary and solar radar measurements, the coordinate distances are much larger than the impact parameter:
\begin{equation}
    z_1, z_2 \gg b.
\end{equation}

In this limit, the logarithmic term simplifies algebraically to the standard weak-field Shapiro form:
\begin{equation}
    \Delta t_{\mathrm{delay}} \approx \frac{2GM}{c_{\mathrm{E}}^3} \ln \left( \frac{4z_1z_2}{b^2} \right).
\end{equation}

This confirms that the localized temporal and spatial impedances encoded by the $2\alpha_{\mathrm{G}}$ refractive correction integrate macroscopically to reproduce the standard weak-field Shapiro signal delay.